% =============================================================================
% Bloque 8: Discusión, Utilidad Real, Cierre (slides 25–27)
% =============================================================================

\section{Discusión y aplicaciones}

% =============================================================================
% [DIVULGACIÓN] Slide oculta para la versión de público general.
% Reactivar quitando \iffalse ... \fi.
% =============================================================================
\iffalse
% --- Slide 25: Contribuciones agrupadas ---
\begin{frame}{Contribuciones de la tesis}
\small
\begin{columns}[T]

\begin{column}{0.33\textwidth}
\setbeamercolor{block title}{fg=white,bg=uaiblue}
\begin{block}{Metodológica}
\scriptsize
\begin{itemize}\setlength\itemsep{0pt}
  \item Un control de calidad para textos de IA, basado en el mapa.
  \item Pesos por calidad, en vez de aceptar/rechazar en seco.
  \item Funciona con cualquier IA, sin tocarla.
\end{itemize}
\end{block}
\end{column}

\begin{column}{0.33\textwidth}
\setbeamercolor{block title}{fg=white,bg=uaiorange}
\begin{block}{Empírica}
\scriptsize
\begin{itemize}\setlength\itemsep{0pt}
  \item +2.25 pp sobre SMOTE\\(comprobado, \textit{p}\,$<$\,0.0001).
  \item +1.10 pp sobre usar la IA sin control de calidad.
  \item Le gana a 9 métodos rivales, clásicos y modernos.
  \item Robusto: 5 IAs y 4 mapas distintos.
\end{itemize}
\end{block}
\end{column}

\begin{column}{0.33\textwidth}
\setbeamercolor{block title}{fg=white,bg=uaigreen}
\begin{block}{Generalización}
\scriptsize
\begin{itemize}\setlength\itemsep{0pt}
  \item Donde más duele, más ayuda: +10.44 pp en clases difíciles.
  \item Otra tarea (NER): +9.61 pp sin modificar nada.
  \item Hallazgo: lo simple le gana a lo complejo.
\end{itemize}
\end{block}
\end{column}

\end{columns}

\vspace{0.5em}
\begin{insight}
\scriptsize\centering
\textbf{Insight central:} la posición en el mapa basta para medir la calidad de un ejemplo\\
— en cualquier tarea — siempre que el criterio sea simple.
\end{insight}
\end{frame}
\fi

% --- Slide 26: Utilidad real / aplicaciones ---
\begin{frame}{¿Cuándo y cómo usarlo? Utilidad práctica}
\small
\begin{columns}[T]

\begin{column}{0.55\textwidth}
\textbf{Receta de aplicación:}
\begin{enumerate}\setlength\itemsep{1pt}
  \item Tienes \alert{10--25 ejemplos por clase}.
  \item Tienes \alert{4--10 clases}.
  \item Pides a una IA el triple de ejemplos.
  \item Filtras por \textbf{cercanía} en el mapa (nivel 1).
  \item Pesas por calidad y entrenas un\\
        clasificador \textbf{lineal}.
\end{enumerate}

\vspace{0.4em}
\setbeamercolor{block title}{fg=white,bg=uaiorange}
\begin{block}{Costo real}
$\sim 10$--$50$ llamadas LLM por clase.\\
A precios actuales: \textbf{centavos por dataset}.
\end{block}
\end{column}

\begin{column}{0.45\textwidth}
\textbf{Posibles casos de aplicación:}
\scriptsize
\begin{itemize}\setlength\itemsep{1pt}
  \item Detección de discurso de odio\\\cita{hate\_speech: +13 pp en clase rara}.
  \item Clasificación de emociones en texto\\\cita{emotion: +4.55 pp; +16.9 pp en \textit{anger}}.
  \item Categorización de noticias por tema\\\cita{20newsgroups, ag\_news: +5 a +15 pp}.
  \item Clasificación de tipo de preguntas\\\cita{trec6: +10.12 pp}.
  \item Detección de spam SMS\\\cita{sms\_spam: $\sim +3$ pp}.
\end{itemize}

\vspace{0.4em}
\small
\textbf{Quién no debería usarlo:}
\scriptsize
\begin{itemize}\setlength\itemsep{1pt}
  \item Más de 50 ejemplos por clase: lo clásico basta.
  \item $\geq 77$ clases: el mapa se fragmenta.
  \item Bosques aleatorios como único clasificador (XGBoost sí sirve).
\end{itemize}
\end{column}

\end{columns}
\end{frame}

% --- Slide 27: Limitaciones, futuro y cierre ---
\begin{frame}{Limitaciones, trabajo futuro y cierre}
\small
\begin{columns}[T]

\begin{column}{0.5\textwidth}
\setbeamercolor{block title}{fg=white,bg=uaiamber}
\begin{block}{Limitaciones honestas}
\scriptsize
\begin{itemize}\setlength\itemsep{0pt}
  \item Evaluación solo en \textbf{inglés}.
  \item Parte de los textos de IA se reutilizó entre repeticiones; verificado aparte: la mejora \textbf{persiste} sin ese atajo.
  \item Pierde efectividad con $\geq 77$ clases.
  \item Las IAs comerciales se niegan a escribir contenido sensible (p.\,ej.\ discurso de odio), justo donde más falta hace.
\end{itemize}
\end{block}
\end{column}

\begin{column}{0.5\textwidth}
\setbeamercolor{block title}{fg=white,bg=uaiblue}
\begin{block}{Trabajo futuro}
\scriptsize
\begin{itemize}\setlength\itemsep{0pt}
  \item Probar en \textbf{español} y otros idiomas.
  \item IAs locales (Llama, Mistral), sin restricciones de contenido.
  \item Que la exigencia de la curva de pesos se ajuste \textbf{sola por clase}.
  \item Otras tareas: responder preguntas, resumir\dots
\end{itemize}
\end{block}
\end{column}

\end{columns}

\vspace{0.6em}
\begin{center}
\setbeamercolor{block title}{fg=white,bg=uaibluedark}
\begin{tcolorbox}[colback=uaibluedark, colframe=uaibluedark, coltext=white,
                  width=0.85\linewidth, halign=center, arc=4pt]
\textbf{Gracias.}\\[0.2em]
{\small Preguntas, comentarios, contraejemplos: bienvenidos.}
\end{tcolorbox}
\end{center}
\end{frame}

% --- Slide: Referencias I ---
\begin{frame}[allowframebreaks]{Referencias}
\scriptsize
\setlength{\parskip}{0.25em}

\textbf{Sobremuestreo geométrico clásico}
\begin{itemize}\setlength\itemsep{0pt}
  \item Chawla, N.\,V., Bowyer, K.\,W., Hall, L.\,O. \& Kegelmeyer, W.\,P. (2002). \textit{SMOTE: Synthetic minority over-sampling technique}. J. Artif. Intell. Res., 16, 321--357.
  \item He, H., Bai, Y., Garcia, E.\,A. \& Li, S. (2008). \textit{ADASYN: Adaptive synthetic sampling approach for imbalanced learning}. IEEE IJCNN, 1322--1328.
\end{itemize}

\textbf{Aumentación textual ligera}
\begin{itemize}\setlength\itemsep{0pt}
  \item Wei, J. \& Zou, K. (2019). \textit{EDA: Easy Data Augmentation Techniques for Boosting Performance on Text Classification Tasks}. EMNLP-IJCNLP, 6382--6388.
  \item Sennrich, R., Haddow, B. \& Birch, A. (2016). \textit{Improving Neural Machine Translation Models with Monolingual Data}. ACL.
  \item Edunov, S., Ott, M., Auli, M. \& Grangier, D. (2018). \textit{Understanding Back-Translation at Scale}. EMNLP, 489--500.
  \item Kobayashi, S. (2018). \textit{Contextual Augmentation: Data Augmentation by Words with Paradigmatic Relations}. NAACL-HLT, 452--457.
  \item Wu, X., Lv, S., Zang, L., Han, J. \& Hu, S. (2019). \textit{Conditional BERT Contextual Augmentation}. ICCS, 84--95.
\end{itemize}

\textbf{Aumentación en espacios de embeddings}
\begin{itemize}\setlength\itemsep{0pt}
  \item Zhang, H., Cisse, M., Dauphin, Y.\,N. \& Lopez-Paz, D. (2018). \textit{mixup: Beyond Empirical Risk Minimization}. ICLR.
  \item Chen, J., Yang, Z. \& Yang, D. (2020). \textit{MixText: Linguistically-Informed Interpolation of Hidden Space for Semi-Supervised Text Classification}. ACL, 2147--2157.
  \item Bayer, M., Kaufhold, M.\,A. \& Reuter, C. (2023). \textit{A Survey on Data Augmentation for Text Classification}. ACM Computing Surveys, 55(7), 1--39.
\end{itemize}

\framebreak

\textbf{Generación con LLMs}
\begin{itemize}\setlength\itemsep{0pt}
  \item Brown, T. et al. (2020). \textit{Language Models Are Few-Shot Learners}. NeurIPS, 33, 1877--1901.
  \item Whitehouse, C., Choudhury, M. \& Aji, A.\,F. (2023). \textit{LLM-powered Data Augmentation for Cross-lingual Commonsense Reasoning}. ACL, 1234--1245.
  \item Suhaeni, C. \& Yong, H.-S. (2023). \textit{Mitigating Class Imbalance in Sentiment Analysis through GPT-3-Generated Synthetic Sentences}. Appl.\ Sci., 13, 9766. ($+12.7\%$ acc.\ en minoritarias.)
  \item Gopali, S., Abri, F., Namin, A.\,S. \& Jones, K.\,S. (2024). \textit{The Applicability of LLMs in Generating Textual Samples for Analysis of Imbalanced Datasets}. IEEE Access, 12, 136451--136465.
  \item Cegin, J., Simko, J. \& Brusilovsky, P. (2025). \textit{LLMs vs Established Text Augmentation Techniques for Classification: When do the Benefits Outweigh the Costs?} NAACL, 10543--10570.
  \item Springer, J.\,M., Devoto, C., Tosheva, M.\,I. \& Ramesh, P. (2025). \textit{On the Influence of Synthetic Data for Few-Shot Text Classification}. Findings ACL, 8249--8268.
\end{itemize}

\textbf{Fine-tuning de LLMs para generación sintética}
\begin{itemize}\setlength\itemsep{0pt}
  \item Cloutier, N.\,A. \& Japkowicz, N. (2023). \textit{Fine-tuned generative LLM oversampling can improve performance over traditional techniques on multiclass imbalanced text classification}. IEEE Big Data, 5181--5186.
  \item Nguyen, D. et al. (2024). \textit{Large Language Models for Imbalanced Classification: Diversity makes the difference} (ImbLLM). arXiv:2510.09783.
  \item Sundararajan, N.\,T. \& Kumarapathirage, A. (2023). \textit{SentiGEN: A Sentiment-Aware Generative Model for Text Data Augmentation}. Findings EMNLP, 9138--9149.
\end{itemize}

\framebreak

\textbf{Filtrado, anomalías y selección de datos}
\begin{itemize}\setlength\itemsep{0pt}
  \item Breunig, M.\,M., Kriegel, H.-P., Ng, R.\,T. \& Sander, J. (2000). \textit{LOF: Identifying Density-Based Local Outliers}. ACM SIGMOD, 93--104.
  \item Robinson, M., Dey, S. \& Chiang, T. (2025). \textit{Token Embeddings Violate the Manifold Hypothesis}. NeurIPS.
  \item Tessari, F. \& Hogan, N. (2024). \textit{Surpassing Cosine Similarity for Multidimensional Comparisons: Dimension Insensitive Euclidean Metric}. arXiv:2407.08623.
  \item Du, Y. et al. (2025). \textit{Disentangling the Roles of Representation and Selection in Data Pruning}. ACL, 16791--16809.
  \item Yu, T., Zhuang, Y., Zhang, Z., Socher, R. \& Xiong, C. (2025). \textit{R.I.P.: Better Models by Survival of the Fittest Data}. ICML.
  \item Liu, Z., Tang, Y., Wang, H. \& Li, M. (2025). \textit{Residual Alignment Model: A Framework for Importance Sampling in Data Augmentation}. NeurIPS.
  \item Tokdar, S.\,T. \& Kass, R.\,E. (2010). \textit{Importance Sampling: A Review}. WIREs Comput.\ Stat., 2(1), 54--60.
\end{itemize}

\textbf{NER y extensiones}
\begin{itemize}\setlength\itemsep{0pt}
  \item Zhou, R. et al. (2022). \textit{MELM: Data Augmentation with Masked Entity Language Modeling for Low-Resource NER}. ACL, 2251--2262.
\end{itemize}

\textbf{Manifold hypothesis y embeddings}
\begin{itemize}\setlength\itemsep{0pt}
  \item Bengio, Y., Courville, A. \& Vincent, P. (2013). \textit{Representation Learning: A Review and New Perspectives}. IEEE TPAMI, 35(8), 1798--1828.
  \item Fefferman, C., Mitter, S. \& Narayanan, H. (2016). \textit{Testing the Manifold Hypothesis}. J. Amer.\ Math.\ Soc., 29(4), 983--1049.
  \item Muennighoff, N., Tazi, N., Magne, L. \& Reimers, N. (2023). \textit{MTEB: Massive Text Embedding Benchmark}. EACL, 2014--2037.
\end{itemize}
\end{frame}
